\chapter{Reading, graded by source quality}
\label{app:reading}

Everything this book cites is in \code{sources.json}, with a URL, the date it
was checked, a quality grade, and whether it was read directly or through
somebody else's coverage. This appendix is the reading list, organised by what
each source is good for.

\section{How the grades work}

\begin{kittable}{@{}lX@{}}
\textbf{Grade} & \textbf{Means} \\
\midrule
\code{strong} & several independent sources, or one with a stated method over a named corpus \\
\code{moderate} & one good source, corroborated in outline elsewhere \\
\code{single-source} & one source; a claim citing it must say so, in the sentence \\
\code{unverified} & traced only to content-marketing; usable only in Appendix~\ref{app:refused} \\
\end{kittable}

A second field records whether the source was read directly or through
coverage. \judgement{That distinction is not decoration. Several primary
sources for this book were unreachable from the machine it was written on, and
a summary of a source is not the source \dash{} it is somebody else's reading
of it, with their emphasis.}

\section{Start here}

\textbf{A field guide built from a scraped corpus of job descriptions,
disclosed processes and take-home repositories.}\source{field-guide} The single
most useful document available for this transition, and the backbone of
Part~V's ordering. One author's methodology over a named corpus, which is more
than anything else in this space offers. Its sub-files on the interview
process, take-homes and the backend learning path are each worth reading
directly.

\textbf{The standard text on AI engineering as a
discipline.}\source{aie-book} Used here for the boundary between this work
and classical machine learning.

\section{On the level}

\textbf{Published engineering ladders.}\source{progression-fyi} Read two or
three rather than one. The most cited of them\source{etsy-ladder} supplies the
word Chapter~\ref{ch:staff} is built on.

\textbf{The archetype taxonomy}\source{staff-archetypes} \textbf{and its two
published rebuttals}\source{archetypes-rebuttal-goedecke}\source{archetypes-rebuttal-ewerlof}.
Read all three. The disagreement is more useful than any of the three
positions.

\textbf{On work that is necessary and invisible to a ladder.}\source{being-glue}
Short, and the clearest statement of why scope rather than effort is what gets
assessed.

\section{On evaluation}

\textbf{Practitioner writing on evaluation and on model-based
grading.}\source{hamel-evals-faq}\source{hamel-llm-judge} The second is the
source for anchored rubrics and calibration in Chapter~\ref{ch:evals}.

\textbf{A model provider's own account of evaluating
agents.}\source{anthropic-agent-evals}

\section{On the changing format}

\textbf{One company's published account of an AI-native
loop}\source{sierra-ai-native-interview} and \textbf{another's statement that
assistance is expected}\source{canva-ai-interviews}. Both were unreachable
directly and are recorded as read through coverage.

\textbf{One industry survey of hiring leaders.}\source{coderpad-hiring-2026} A
vendor's own survey, used only for the coarse split.

\section{On safety}

\textbf{On injection in tool-using systems and blast
radius}\source{mcp-prompt-injection} \textbf{and on why agency makes it
worse}\source{agentic-injection}.

\section{The register, in full}

Generated from \code{sources.json} rather than typed, so the list and the
numbers in the margins of every claim cannot come apart. The number here is
the number in the superscript.

\kitsourcetable

\section{What is deliberately not here}

No preparation-vendor question lists. No aggregated ``top fifty questions''
pages. Appendix~\ref{app:refused} explains why: that genre recycles itself, and
several sources repeating one origin produce the appearance of consensus from a
single post.

\judgement{This list is seeded rather than complete, and it is the thinnest
part of the book. The repository's own register says so, which is the habit
Chapter~\ref{ch:ledger} is about, applied to the thing you are holding.}
